\begin{table}[htbp]
\centering
\small
\caption{Problemas B e C, Passo 3 (openai/gpt-4o-mini): transferência zero-shot contra cada métrica estimada no A.}
\label{tab:BC_passo3_todas_metricas_A_gpt4mini}
\begin{tabular}{llllrrr}
\toprule
Prob. & Métrica de origem (A) & Algoritmo & Consistência (\%) & $\kappa$ & F1-Macro (\%) & Acur.\ LLM vs real \\
\midrule
B & zero-shot & perceptron & 69,0 (67,0–73,0) & 0,320 (0,255–0,414) & 64,8 (61,8–70,1) & 56,3 (56,0–57,0) \\
B & zero-shot & nnls & 70,3 (67,0–75,0) & 0,346 (0,291–0,446) & 66,1 (62,4–71,9) & 56,3 (56,0–57,0) \\
B & 4-shot & perceptron & 69,8 (54,0–78,0) & 0,339 (0,060–0,508) & 65,7 (50,1–75,1) & 56,3 (56,0–57,0) \\
B & 4-shot & nnls & 67,4 (48,0–80,0) & 0,300 (-0,077–0,543) & 63,4 (42,1–76,6) & 56,3 (56,0–57,0) \\
B & 10-shot & perceptron & 74,3 (69,0–80,0) & 0,422 (0,316–0,520) & 70,3 (64,1–75,7) & 56,3 (56,0–57,0) \\
B & 10-shot & nnls & 71,9 (48,0–80,0) & 0,378 (-0,077–0,532) & 67,8 (42,1–76,2) & 56,3 (56,0–57,0) \\
B & 20-shot & perceptron & 73,2 (70,0–80,0) & 0,400 (0,330–0,520) & 69,1 (64,9–75,7) & 56,3 (56,0–57,0) \\
B & 20-shot & nnls & 75,8 (70,0–81,0) & 0,452 (0,330–0,543) & 71,9 (64,9–76,7) & 56,3 (56,0–57,0) \\
\textbf{B} & \textbf{40-shot} & \textbf{perceptron} & \textbf{75,3 (70,0–80,0)} & \textbf{0,441 (0,330–0,520)} & \textbf{71,3 (64,9–75,7)} & \textbf{56,3 (56,0–57,0)} \\
\textbf{B} & \textbf{40-shot} & \textbf{nnls} & \textbf{75,0 (70,0–81,0)} & \textbf{0,437 (0,330–0,561)} & \textbf{71,0 (64,9–77,6)} & \textbf{56,3 (56,0–57,0)} \\
B & W ajustado ao rótulo real (teto) & perceptron & 72,0 (70,0–76,0) & 0,381 (0,330–0,484) & 68,0 (64,9–73,6) & 56,3 (56,0–57,0) \\
B & W ajustado ao rótulo real (teto) & nnls & 77,0 (75,0–80,0) & 0,469 (0,403–0,520) & 72,9 (69,3–75,7) & 56,3 (56,0–57,0) \\
\midrule
C & zero-shot & perceptron & 87,0 (83,0–90,0) & 0,725 (0,636–0,789) & 86,1 (81,4–89,4) & 18,0 (16,0–21,0) \\
C & zero-shot & nnls & 87,7 (84,0–90,0) & 0,739 (0,654–0,792) & 86,8 (82,4–89,5) & 18,0 (16,0–21,0) \\
C & 4-shot & perceptron & 85,9 (81,0–89,0) & 0,707 (0,613–0,771) & 85,1 (80,3–88,4) & 18,0 (16,0–21,0) \\
C & 4-shot & nnls & 83,6 (78,0–88,0) & 0,659 (0,551–0,752) & 82,6 (77,1–87,4) & 18,0 (16,0–21,0) \\
C & 10-shot & perceptron & 86,1 (84,0–88,0) & 0,711 (0,654–0,754) & 85,3 (82,4–87,5) & 18,0 (16,0–21,0) \\
C & 10-shot & nnls & 85,0 (78,0–88,0) & 0,689 (0,551–0,752) & 84,2 (77,1–87,4) & 18,0 (16,0–21,0) \\
C & 20-shot & perceptron & 86,2 (84,0–88,0) & 0,713 (0,654–0,752) & 85,4 (82,4–87,4) & 18,0 (16,0–21,0) \\
C & 20-shot & nnls & 85,8 (84,0–88,0) & 0,705 (0,660–0,752) & 85,0 (82,6–87,4) & 18,0 (16,0–21,0) \\
\textbf{C} & \textbf{40-shot} & \textbf{perceptron} & \textbf{86,1 (84,0–88,0)} & \textbf{0,711 (0,660–0,752)} & \textbf{85,3 (82,6–87,4)} & \textbf{18,0 (16,0–21,0)} \\
\textbf{C} & \textbf{40-shot} & \textbf{nnls} & \textbf{85,9 (84,0–88,0)} & \textbf{0,707 (0,660–0,752)} & \textbf{85,1 (82,6–87,4)} & \textbf{18,0 (16,0–21,0)} \\
C & W ajustado ao rótulo real (teto) & perceptron & 86,0 (84,0–87,0) & 0,709 (0,660–0,735) & 85,2 (82,6–86,5) & 18,0 (16,0–21,0) \\
C & W ajustado ao rótulo real (teto) & nnls & 85,0 (84,0–86,0) & 0,690 (0,660–0,713) & 84,2 (82,6–85,4) & 18,0 (16,0–21,0) \\
\bottomrule
\multicolumn{7}{p{0.95\linewidth}}{\footnotesize Rótulos ZERO-SHOT do LLM em B e C (100 pontos novos cada); só $\hat{W}$ é transferido do A (R$^2$), os centróides são recalculados com os rótulos do LLM no próprio problema (consistência). $\kappa$ = Cohen's $\kappa$ (faixas de Landis e Koch, 1977). F1-Macro entre a métrica e o LLM. Acur.\ LLM vs real = acurácia (vs rótulo real) do LLM, contexto. Em zero-shot (e com nomes de classe neutros) a correspondência entre nome e classe é arbitrária: a acurácia (vs rótulo real) do LLM pode sair abaixo de 50\%; o valor invariante à permutação é max($a$, 1 $-$ $a$). Em negrito: as linhas (Perceptron e NNLS) do regime/k eleito pelo runner (coluna eleita\_passo6: maior fidelidade média no A, Perceptron, R$^2$, zero ou few-shot), excluído o teto. Passos de B e C: o LLM rotula os 100 pontos em zero-shot (Passo 2) e cada métrica do A é aplicada a esses rótulos (Passo 3). A métrica de origem vem do Passo 4 do A (zero-shot = rótulos do Passo 2; $k$-shot = rótulos do Passo 3; ``W ajustado ao rótulo real (teto)'' = Passo 5 do A, que não concorre à eleição).} \\
\end{tabular}
\end{table}
